\problemname{Just Take The Stairs!}
%\illustration{0.3}{image.jpg}{
%    Caption of the illustration (optional).
%    CC BY-SA 4.0 by X on \href{https://example.com/reference-to-image}{Y}
%}

It's 7{:}55 Monday morning at the IT-University of Copenhagen (ITU). A massive stream of students
is packed into the beautiful ITU Atrium on the ground floor (floor $0$), bleary-eyed and clutching
laptops. Analog hasn't opened yet, so nobody has had coffee, and tempers are short.

The stairs are right there, next to the elevators. Walking up two floors would take fifteen seconds.
But without coffee in their systems, no one can be expected to climb stairs, and so every last
student shuffles into the elevator queue anyway.

To make matters worse, ITU just suffered severe water damage, and one of the two elevators has
broken! (The other side of the building has two fully functional ones, but nobody is walking that
far before coffee.)

As the Facility Management Intern, you have been given the job of managing the last remaining
elevator and minimizing the wear and tear on it. (You cannot force anyone to take the stairs
they all have noise-cancelling headphones on and simply cannot hear you.) Your solution? Minimize
the total work the elevator has to do.

The elevator has a capacity of $C$ persons and serves floors $0$ (ground) through $5$. It starts
at floor $0$.

There are $N$ groups of students in the queue. Each group consists of $S_i$ students all headed for
the same floor $F_i$. Danish queue culture is strict: groups must be served in the exact order they
stand, and you may not split a group across trips.

A trip carries one or more \emph{consecutive} groups, with total persons $\leq C$.

\subsection*{Wear and Tear}

Each trip of the elevator incurs the following cost:

\begin{enumerate}
    \item \textbf{Travel up:} $1$ per floor, from floor $0$ up to the highest destination floor
    among the groups in the trip. Passing intermediate floors along the way adds no extra cost.
    \item \textbf{Stop cost:} 1 for each distinct destination floor among the groups in the trip.
    Boarding at floor 0 is free; only disembarking stops are counted. For example, a trip carrying
    3 groups bound for floors $\{2, 5, 5\}$ has a stop cost of $2$, since there are two distinct
    destination floors.
    \item \textbf{Travel down:} $1$ per floor, back to floor $0$ to pick up the next group. After
    the \emph{final} trip, the elevator does not need to return to floor $0$, so no descent cost
    is incurred for the final trip.
\end{enumerate}

\begin{Input}
    The input consists of:
    \begin{itemize}
        \item One line containing two integers $N$ ($1 \leq N \leq 200\,000$) and
        $C$ ($1 \leq C \leq 50$).
        \item The next $N$ lines each contain two integers which defines a group: $S_i$ ($1 \leq S_i \leq C$), the amount of students and $F_i$ ($1 \leq F_i \leq 5$), the floor that group must get off at.
    \end{itemize}
\end{Input}

\begin{Output}
    Output a single integer: the minimum total cost.
\end{Output}

\subsection*{Explanation of Sample Input 1}
\textbf{Trip 1:} Group 1 only (1 up + 1 open + 1 down) = 3 \\
\textbf{Trip 2:} Group 2 \& 3 (5 up + 1 open) = 6 \\
\textbf{Total:} 9
